\documentclass[aps,preprint,onecolumn,nofootinbib,superscriptaddress]{revtex4-2}

\usepackage{amsmath,amssymb,amsthm,mathtools,bm}
\usepackage[colorlinks=true,linkcolor=blue,citecolor=blue,urlcolor=blue]{hyperref}
\usepackage{microtype}

\newcommand{\I}{\mathbf{1}}
\newcommand{\E}{\mathbb{E}}
\newcommand{\Prob}{\mathbb{P}}
\newcommand{\Tr}{\operatorname{Tr}}
\newcommand{\supp}{\operatorname{supp}}
\newcommand{\dd}{\mathrm{d}}
\newcommand{\cA}{\mathcal{A}}
\newcommand{\cB}{\mathcal{B}}
\newcommand{\cK}{\mathcal{K}}
\newcommand{\cL}{\mathcal{L}}
\newcommand{\cM}{\mathcal{M}}
\newcommand{\cP}{\mathcal{P}}
\newcommand{\cX}{\mathcal{X}}
\newcommand{\PF}{\mathfrak{P}}
\newcommand{\Koop}{\mathfrak{K}}

\theoremstyle{definition}
\newtheorem{definition}{Definition}
\newtheorem{remark}{Remark}

\theoremstyle{plain}
\newtheorem{proposition}{Proposition}

\begin{document}

\title{Markov Transfer Operators for Covariance-Matrix Dynamics}
\author{}
\date{\today}

\begin{abstract}
This note records a precise formulation of the stochastic Gaussian
covariance update as a discrete-time Markov process on covariance matrices.
Each elementary circuit event is a state-dependent random rational map
\(G\mapsto F_\alpha(G)\), where the branch label \(\alpha\) contains the local
operation and measurement outcome. The corresponding Markov kernel pushes
forward probability distributions on the covariance manifold, while its
adjoint acts on observables by the usual Koopman transfer operator. Linearizing
the same random maps gives a multiplicative tangent cocycle
\(\delta G_{t+1}=T_{\alpha_t}(G_t)^\dagger\delta G_t T_{\alpha_t}(G_t)\).
The Lyapunov spectrum is the quenched singular-value spectrum of this cocycle,
not the spectrum of an outcome-averaged transfer matrix.
\end{abstract}

\maketitle

\section{Covariance state space}

Let \(G\in\mathbb{C}^{N\times N}\) denote the complex covariance matrix of the
top-layer Gaussian state. In the convention used here,
\begin{equation}
  G_{ij}=2\langle c_i^\dagger c_j\rangle-\delta_{ij}.
\end{equation}
Thus the one-body occupation matrix is
\begin{equation}
  C=\frac{1}{2}(G+\I).
\end{equation}
The natural state space is a real manifold, or a stratified real analytic set,
of admissible covariance matrices,
\begin{equation}
  \cX=\{G:\,G=G^\dagger,\,-\I\leq G\leq \I,\text{ plus the imposed symmetry
  constraints}\}.
\end{equation}
Pure Gaussian states lie on the boundary component satisfying
\begin{equation}
  G^2=\I.
\end{equation}
A perturbation away from a pure state obeys
\begin{equation}
  (G+\delta G)^2
  =
  \I+\{G,\delta G\}+O(\delta G^2).
\end{equation}
Therefore a first-order tangent perturbation along the pure-state manifold
satisfies \(\{G,\delta G\}=0\), whereas first-order mixedness is encoded by the
normal component for which \(\{G,\delta G\}\neq0\).

\section{Elementary branch maps}

An elementary realized Gaussian channel branch acts on \(G\) by a fractional
linear covariance map
\begin{equation}
  F_\alpha(G)
  =
  C_\alpha
  -
  B_\alpha^\dagger
  \left(A_\alpha-G^{-1}\right)^{-1}
  B_\alpha,
  \label{eq:branch-map}
\end{equation}
where \(\alpha\) labels the operation, site or mode, direction, and measurement
outcome. The domain of this formula is the open subset of \(\cX\) on which the
inverse exists. Singular events should be understood either by restricting to
the nonsingular support block or by taking the appropriate Schur-complement
limit.

For a filling or deletion operation at mode \((r,\nu)\), a useful branch label is
\begin{equation}
  \alpha=(r,\nu,\sigma,s),
  \qquad
  \sigma\in\{\mathrm{fill},\mathrm{del}\},
  \qquad
  s\in\{0,1\}.
\end{equation}
Let
\begin{equation}
  P_{r,\nu}=W_{r,\nu}W_{r,\nu}^\dagger,
  \qquad
  Q_{r,\nu}=\I-P_{r,\nu}.
\end{equation}
The branch probability is state dependent. For instance, for an occupation
measurement of the mode \(W_{r,\nu}\),
\begin{equation}
  p_{r,\nu}(s=1\mid G)
  =
  \Tr\!\left(P_{r,\nu}\frac{G+\I}{2}\right),
  \qquad
  p_{r,\nu}(s=0\mid G)
  =
  1-p_{r,\nu}(s=1\mid G),
  \label{eq:born-probability}
\end{equation}
up to the precise sign convention used to associate the two branch labels with
occupied and empty outcomes. More general schedules may also include an
externally chosen probability for selecting \((r,\nu,\sigma)\).

\section{Discrete-time Markov process on covariance matrices}

The covariance process is a time-inhomogeneous Markov chain on \(\cX\). Given
the current covariance \(G_t\), one samples a branch label \(\alpha_t\) from a
conditional probability law
\begin{equation}
  \Prob(\alpha_t=\alpha\mid G_t=G)=p_t(\alpha\mid G),
\end{equation}
and then updates deterministically inside that branch,
\begin{equation}
  G_{t+1}=F_\alpha(G_t).
  \label{eq:discrete-process}
\end{equation}
Equivalently,
\begin{equation}
  G_{t+1}=F_{\alpha_t}(G_t),
  \qquad
  \alpha_t\sim p_t(\cdot\mid G_t).
\end{equation}
The Markov property is immediate:
\begin{equation}
  \Prob(G_{t+1}\in E\mid G_t,G_{t-1},\ldots,G_0)
  =
  \Prob(G_{t+1}\in E\mid G_t)
\end{equation}
for every Borel set \(E\subset\cX\), provided the circuit schedule at time \(t\)
is fixed or is itself included as part of the time-dependent law \(p_t\).

The associated Markov kernel is
\begin{equation}
  K_t(G,E)
  =
  \sum_{\alpha\in\cA_t}
  p_t(\alpha\mid G)\,
  \bm{1}_E\!\left(F_\alpha(G)\right),
  \label{eq:markov-kernel}
\end{equation}
or, in distributional notation,
\begin{equation}
  K_t(G,\dd G')
  =
  \sum_{\alpha\in\cA_t}
  p_t(\alpha\mid G)\,
  \delta_{F_\alpha(G)}(\dd G').
  \label{eq:kernel-delta}
\end{equation}
Here \(\cA_t\) is the finite set of branch labels available at step \(t\). For a
deterministic choice of location and operation, \(\cA_t\) contains only the two
measurement outcomes. For a randomly selected location or operation, \(\cA_t\)
contains the combined operation-outcome labels and \(p_t\) includes both the
external sampling probability and the Born probability.

\section{Transfer operators on observables and distributions}

Let \(\varphi:\cX\to\mathbb{C}\) be a bounded measurable observable of the
covariance matrix. The Koopman, or backward, Markov transfer operator is
\begin{equation}
  (\Koop_t\varphi)(G)
  =
  \E\!\left[\varphi(G_{t+1})\mid G_t=G\right]
  =
  \sum_{\alpha\in\cA_t}
  p_t(\alpha\mid G)\,
  \varphi\!\left(F_\alpha(G)\right).
  \label{eq:koopman}
\end{equation}
The dual Perron-Frobenius operator evolves probability measures. If \(\mu_t\)
is the law of \(G_t\), then
\begin{equation}
  \mu_{t+1}=\PF_t\mu_t,
  \qquad
  (\PF_t\mu_t)(E)
  =
  \int_\cX K_t(G,E)\,\mu_t(\dd G).
  \label{eq:measure-pushforward}
\end{equation}
Equivalently, duality with observables gives
\begin{equation}
  \int_\cX \varphi(G')\,\mu_{t+1}(\dd G')
  =
  \int_\cX (\Koop_t\varphi)(G)\,\mu_t(\dd G).
  \label{eq:duality}
\end{equation}

If \(\mu_t\) has a density \(\rho_t\) with respect to a chosen reference measure
\(\dd G\) on the real covariance manifold, then formally
\begin{equation}
  \rho_{t+1}(G')
  =
  \sum_{\alpha\in\cA_t}
  \int_\cX
  \rho_t(G)\,
  p_t(\alpha\mid G)\,
  \delta\!\left(G'-F_\alpha(G)\right)
  \dd G.
  \label{eq:density-delta}
\end{equation}
When each branch map is locally invertible on the relevant smooth stratum, this
can be written as
\begin{equation}
  \rho_{t+1}(G')
  =
  \sum_{\alpha\in\cA_t}
  \sum_{G\in F_\alpha^{-1}(G')}
  \frac{\rho_t(G)\,p_t(\alpha\mid G)}
       {|\det_{\mathbb{R}} D F_\alpha(G)|}.
  \label{eq:pf-density}
\end{equation}
The determinant in Eq.~\eqref{eq:pf-density} is the real Jacobian on the
appropriate tangent space of admissible covariance matrices. In practice,
Eq.~\eqref{eq:measure-pushforward} is the more robust definition: measurement
maps can be rank deficient or singular on lower-dimensional sets, in which case
the evolved measure need not remain absolutely continuous.

\section{Full layers and unit cells}

A full circuit layer or unit cell consists of a finite sequence of elementary
updates. Suppose the physical order is
\begin{equation}
  j=1,2,\ldots,m,
\end{equation}
with elementary kernels \(K_t^{(j)}\). The covariance process inside the cell is
\begin{equation}
  G^{(j)}=F_{\alpha^{(j)}}^{(j)}(G^{(j-1)}),
  \qquad
  \alpha^{(j)}\sim p_t^{(j)}(\cdot\mid G^{(j-1)}).
\end{equation}
The cell-level kernel is the kernel composition
\begin{equation}
  K_t^{\mathrm{cell}}
  =
  K_t^{(1)}K_t^{(2)}\cdots K_t^{(m)}
\end{equation}
in the following explicit sense:
\begin{align}
  K_t^{\mathrm{cell}}(G,E)
  &=
  \int_{\cX}\cdots\int_{\cX}
  K_t^{(1)}(G,\dd G^{(1)})
  K_t^{(2)}(G^{(1)},\dd G^{(2)})
  \cdots
  K_t^{(m)}(G^{(m-1)},E).
  \label{eq:kernel-composition}
\end{align}
Thus for observables,
\begin{equation}
  \Koop_t^{\mathrm{cell}}\varphi
  =
  \Koop_t^{(1)}
  \Koop_t^{(2)}
  \cdots
  \Koop_t^{(m)}
  \varphi,
  \label{eq:observable-composition}
\end{equation}
where the first physical operation appears first because
\((\Koop_t^{(1)}\psi)(G)=\E[\psi(G^{(1)})\mid G^{(0)}=G]\). For measures the
dual order is
\begin{equation}
  \mu_t^{\mathrm{out}}
  =
  \PF_t^{(m)}
  \cdots
  \PF_t^{(2)}
  \PF_t^{(1)}
  \mu_t^{\mathrm{in}}.
  \label{eq:measure-composition}
\end{equation}

\section{Linearization and tangent transfer matrices}

The branchwise tangent map follows by differentiating
Eq.~\eqref{eq:branch-map}. Let
\begin{equation}
  M_\alpha(G)=A_\alpha-G^{-1}.
\end{equation}
For a perturbation \(\delta G\),
\begin{equation}
  \delta G^{-1}
  =
  -G^{-1}\delta G\,G^{-1},
\end{equation}
and therefore
\begin{equation}
  \delta M_\alpha
  =
  G^{-1}\delta G\,G^{-1}.
\end{equation}
Since \(\delta M^{-1}=-M^{-1}(\delta M)M^{-1}\), one obtains
\begin{align}
  D F_\alpha(G)[\delta G]
  &=
  B_\alpha^\dagger
  M_\alpha(G)^{-1}
  G^{-1}\delta G\,G^{-1}
  M_\alpha(G)^{-1}
  B_\alpha .
\end{align}
Using the noncommutative identities
\begin{equation}
  M_\alpha(G)^{-1}G^{-1}
  =
  (G A_\alpha-\I)^{-1},
  \qquad
  G^{-1}M_\alpha(G)^{-1}
  =
  (A_\alpha G-\I)^{-1},
\end{equation}
this becomes
\begin{equation}
  D F_\alpha(G)[\delta G]
  =
  B_\alpha^\dagger
  (G A_\alpha-\I)^{-1}
  \delta G
  (A_\alpha G-\I)^{-1}
  B_\alpha.
  \label{eq:linearization}
\end{equation}
When the covariance and Choi blocks obey the Hermiticity relation needed to
identify the left factor as the adjoint of the right factor, define
\begin{equation}
  T_\alpha(G)
  =
  (A_\alpha G-\I)^{-1}B_\alpha.
  \label{eq:tangent-transfer}
\end{equation}
Then
\begin{equation}
  \delta G_{t+1}
  =
  T_{\alpha_t}(G_t)^\dagger
  \delta G_t
  T_{\alpha_t}(G_t).
  \label{eq:tangent-update}
\end{equation}

\section{Lifted Markov process and Lyapunov cocycle}

Equations \eqref{eq:discrete-process} and \eqref{eq:tangent-update} define a
Markov process lifted to the tangent bundle:
\begin{align}
  G_{t+1}
  &=
  F_{\alpha_t}(G_t),
  \\
  \delta G_{t+1}
  &=
  T_{\alpha_t}(G_t)^\dagger
  \delta G_t
  T_{\alpha_t}(G_t),
  \\
  \alpha_t
  &\sim
  p_t(\cdot\mid G_t).
\end{align}
Along one realized trajectory, set
\begin{equation}
  T_t:=T_{\alpha_t}(G_{t-1}).
\end{equation}
Then the chronological product convention compatible with
\(\delta G_t=T_t^\dagger\delta G_{t-1}T_t\) is
\begin{equation}
  \widetilde T_n
  =
  T_1T_2\cdots T_n,
  \label{eq:chronological-product}
\end{equation}
so that
\begin{equation}
  \delta G_n
  =
  \widetilde T_n^\dagger
  \delta G_0
  \widetilde T_n.
  \label{eq:delta-product}
\end{equation}
The Lyapunov exponents are the quenched singular-value growth rates
\begin{equation}
  \lambda_i
  =
  \lim_{n\to\infty}
  \frac{1}{n}\log \sigma_i(\widetilde T_n),
  \label{eq:lyapunov-svd}
\end{equation}
when the limit exists. Equivalently,
\begin{equation}
  2\lambda_i
  =
  \lim_{n\to\infty}
  \frac{1}{n}
  \log \operatorname{eig}_i
  \left(\widetilde T_n^\dagger\widetilde T_n\right).
  \label{eq:lyapunov-positive}
\end{equation}
In a stationary ergodic setting, existence and nonrandomness of the limiting
exponents are the content of the multiplicative ergodic theorem applied to the
Markov cocycle \(T_\alpha(G)\).

\section{Branch averaging versus quenched products}

A useful distinction is between the annealed distributional dynamics of \(G\)
and the quenched tangent dynamics along a realized trajectory. The distribution
of \(G\) evolves by the Markov kernel,
\begin{equation}
  \mu_{t+1}=\PF_t\mu_t.
\end{equation}
However, the tangent update is quadratic in the branch transfer matrix. Thus,
conditional on \(G_t=G\),
\begin{equation}
  \E[\delta G_{t+1}\mid G_t=G,\delta G_t]
  =
  \sum_{\alpha\in\cA_t}
  p_t(\alpha\mid G)\,
  T_\alpha(G)^\dagger
  \delta G_t
  T_\alpha(G).
  \label{eq:average-quadratic}
\end{equation}
This is not equal in general to
\begin{equation}
  \overline T(G)^\dagger\delta G_t\,\overline T(G),
  \qquad
  \overline T(G)=
  \sum_{\alpha\in\cA_t}p_t(\alpha\mid G)T_\alpha(G).
\end{equation}
Consequently the Lyapunov spectrum relevant to relaxation of perturbations is
the spectrum of the random product in Eq.~\eqref{eq:chronological-product}, not
the spectrum of an outcome-averaged matrix.

\section{Summary}

The covariance dynamics can be organized at three levels.
First, the state \(G_t\) is a discrete-time Markov chain with transition kernel
\(K_t(G,\dd G')\). Second, distributions of \(G_t\) evolve by the
Perron-Frobenius pushforward \(\PF_t\), while observables evolve by the dual
Koopman operator \(\Koop_t\). Third, infinitesimal perturbations define a
matrix cocycle over the Markov chain,
\begin{equation}
  (G_t,\delta G_t)
  \mapsto
  \left(F_{\alpha_t}(G_t),
  T_{\alpha_t}(G_t)^\dagger\delta G_tT_{\alpha_t}(G_t)\right).
\end{equation}
The Lyapunov spectrum is obtained from the singular values of the chronological
product of the realized tangent transfer matrices.

\end{document}
